\documentclass[../main.tex]{subfiles}

\begin{document}

\subsection{Overview}

CVE-2023-39191 is a \textbf{verifier logic bug} affecting the handling of eBPF dynamic pointers, or \texttt{dynptr}. The vulnerability affects Linux kernels in the 5.19--6.1.x range and originates from \textbf{incomplete validation of partially overlapping \texttt{dynptr} objects} on the BPF stack. Since BPF programs execute in kernel context after verifier approval, the bug can be transformed from a static-analysis mistake into an \textbf{Out-of-Bounds (OOB) read/write primitive} and, with additional memory shaping, into a full \textbf{Local Privilege Escalation (LPE)}.

\begin{table}[H]
\centering
\begin{tabularx}{\textwidth}{|l|X|}
\hline
\textbf{CVE ID} & CVE-2023-39191 \\
\hline
\textbf{Affected versions} & Linux kernel 5.19--6.1.x \\
\hline
\textbf{Component} & eBPF verifier, dynptr stack-state tracking \\
\hline
\textbf{Vulnerability type} & Type confusion / stale verifier state \\
\hline
\textbf{Root cause} & Missing invalidation of a dynptr after partial stack overlap \\
\hline
\textbf{Impact} & OOB read/write, KASLR bypass, Local Privilege Escalation \\
\hline
\textbf{CVSS} & 8.2 HIGH \\
\hline
\end{tabularx}
\caption{Summary of CVE-2023-39191.}
\end{table}

A \texttt{dynptr} occupies 16 bytes on the BPF stack. The verifier models it as two adjacent 8-byte slots: the first one contains the data pointer, while the second one stores metadata such as size, type, flags and offset. The security invariant is simple: if any byte of a tracked \texttt{dynptr} is overwritten, \textbf{the whole object must be invalidated}. Vulnerable kernels fail to enforce this invariant when a second \texttt{dynptr} is initialized at an offset that overlaps only the upper half of the first one.

The practical exploit developed for this vulnerability does not stop at a verifier bypass. It uses the corrupted \texttt{dynptr} to obtain a stable OOB read/write primitive, calibrates physical memory placement, sprays \texttt{struct cred} objects, scans for UID/GID signatures, overwrites credentials and finally creates a SUID-root shell payload.

\subsection{Vulnerability Analysis}

The vulnerable sequence creates a first local \texttt{dynptr} at \texttt{fp-32}, then creates a second ring-buffer \texttt{dynptr} at \texttt{fp-24}. Since both objects are 16 bytes wide, the second object overwrites the metadata slot of the first one.

\begin{lstlisting}[caption={Overlapping dynptr layout used by the tested exploit.}]
dynptr1 at fp-32:
  [fp-32, fp-24)   data pointer
  [fp-24, fp-16)   size/type/offset metadata

dynptr2 at fp-24:
  [fp-24, fp-16)   data pointer
  [fp-16, fp-8)    size/type/offset metadata

overlap:
  [fp-24, fp-16)   dynptr1 metadata overwritten by dynptr2 data pointer
\end{lstlisting}

At verification time, \texttt{dynptr1} is still considered initialized and safe. At runtime, however, the helper creating \texttt{dynptr2} writes a ring-buffer pointer exactly where \texttt{dynptr1} stores its size and offset metadata. The result is a \textbf{type confusion condition}: the verifier believes it is dealing with a bounded local object, while the kernel executes operations on a corrupted object whose metadata was partially replaced by a kernel address.

The minimal trigger can be summarized as follows:

\begin{lstlisting}[language=C, caption={Core trigger sequence.}]
/* dynptr1: local dynptr backed by a BPF map value */
r1 = map_value;
r2 = 8;
r3 = 0;
r4 = fp - 32;
call bpf_dynptr_from_mem;

/* scrub the upper slot of dynptr1 */
*(u64 *)(fp - 24) = 0;

/* dynptr2: ringbuf dynptr overlapping dynptr1 metadata */
r1 = ringbuf;
r2 = 8;
r3 = 0;
r4 = fp - 24;
call bpf_ringbuf_reserve_dynptr;
\end{lstlisting}

The critical point is not the helper itself, but the verifier state that survives after the overlap. A correct verifier should destroy the old \texttt{dynptr} as soon as the new helper writes into any byte of its 16-byte range. In the vulnerable implementation, stale stack metadata remains trusted long enough to use the corrupted object.

This creates the central verifier/runtime divergence:

\begin{itemize}
    \item \textbf{Verifier view}: \texttt{dynptr1} is still a valid local dynptr backed by a bounded BPF map value.
    \item \textbf{Runtime view}: the metadata of \texttt{dynptr1} has been overwritten by the ring-buffer dynptr pointer.
\end{itemize}

Once this mismatch exists, later \texttt{dynptr} helper calls operate on an object whose logical bounds no longer correspond to the verifier-approved model.

\subsection{Exploit Primitive}

The exploit turns the stale verifier state into an \textbf{OOB memory primitive}. The first \texttt{dynptr} still points to a controlled BPF map value, but its size and offset fields are overwritten by bytes coming from the \textbf{vmalloc-backed ring-buffer pointer}. Because vmalloc addresses contain high bits such as \texttt{0xffffc900}, those bytes can be interpreted as a very large logical access window. In the tested exploit, this produces an effective window of roughly 4 GiB beyond the original map value.

\begin{lstlisting}[caption={Verifier/runtime divergence.}]
Verifier view:
  dynptr1 = local dynptr, bounded to map_value

Runtime view:
  dynptr1 = local data pointer + corrupted vmalloc metadata

Effect:
  map_value + huge corrupted offset -> OOB kernel access
\end{lstlisting}

The BPF payload exposes this primitive to userspace through a BPF map used as a command channel. Userspace writes the desired offset, the value to write and the operation mode; the BPF program performs the corrupted dynptr read/write; userspace then retrieves the result from the same map. Conceptually, the exploit reduces the bug to the following interface:

\begin{lstlisting}[language=C, caption={Abstract exploit interface.}]
oob_read(offset, &value);
oob_write(offset, value);
\end{lstlisting}

This is the key transition: the vulnerability is no longer only a verifier bypass. It becomes a \textbf{reusable kernel read/write primitive} controlled by the userspace orchestrator.

The BPF payload is carefully structured to keep the corrupted object usable while avoiding unnecessary crashes. It creates the overlapped dynptrs, performs the read or write through the corrupted object, stores the result into a map slot and discards the ring-buffer dynptr before exit. This makes the primitive repeatable and stable enough to support the later credential attack.

\subsection{Exploit Architecture}

The exploit is composed of two cooperating layers:

\begin{enumerate}
    \item a userspace C orchestrator;
    \item a BPF payload executed inside the kernel.
\end{enumerate}

The userspace side is responsible for preparing memory, creating the required maps and ring buffers, loading the BPF program, triggering execution and interpreting the leaked data. The BPF payload is responsible for creating the overlapping dynptr layout and performing the actual OOB read/write through the corrupted object.

The exploit uses a small BPF array map as a communication channel. The map stores:

\begin{itemize}
    \item the requested OOB offset;
    \item the value to write;
    \item the result of the last OOB read;
    \item a flag selecting read-only or read/write mode.
\end{itemize}

The high-level BPF execution flow is:

\begin{lstlisting}[caption={High-level BPF payload flow.}]
1. Read command parameters from the BPF map.
2. Create dynptr1 at fp-32.
3. Create dynptr2 at fp-24, overlapping dynptr1.
4. Use corrupted dynptr1 for OOB read.
5. Optionally use corrupted dynptr1 for OOB write.
6. Store the read result back into the map.
7. Discard the ringbuf dynptr and exit.
\end{lstlisting}

The userspace side repeatedly invokes this payload to implement higher-level exploit phases. In practice, every memory probe or write performed during the LPE chain is reduced to one BPF-triggered OOB operation.

\subsection{From OOB Primitive to Root}

The full exploit must solve a practical placement problem: the OOB window points to a physical region determined by the placement of the BPF map value and by the corrupted offset. Therefore, before targeting credentials, the exploit must \textbf{discover and shape where the OOB primitive lands}.

\subsubsection{Adaptive Memory Calibration}

The exploit first performs adaptive memory calibration. After establishing the OOB read/write primitive, it allocates anonymous \texttt{mmap} memory in controlled chunks and fills each allocation with the recognizable pattern \texttt{0xAAAAAAAAAAAAAAAA}. After every allocation step, the exploit probes the OOB window. When the OOB read returns the magic pattern, the exploit knows that the mmap pressure has reached the physical region targeted by the corrupted dynptr offset.

\begin{lstlisting}[caption={Adaptive calibration output.}]
[*] Coarse pass (64MB steps)...
  [ 64] +64MB = 4096MB total *** MAGIC DETECTED ***
[*] Fine pass (2MB steps from 4032MB)...
  [ 64] +2MB = 4034MB total *** MAGIC DETECTED ***
[+] Calibrated: 4034MB pressure
\end{lstlisting}

The calibration is performed in two stages. The coarse phase uses 64 MiB chunks to quickly identify the approximate point where the mmap pressure reaches the OOB window. The exploit then releases the last coarse chunk and repeats the search with 2 MiB chunks, producing a finer placement. In the documented execution, the final calibrated pressure is around 4034 MiB.

In the Buildroot/QEMU test environment, after calibration the mmap pressure covers physical memory up to approximately the OOB target. The first free pages after the mmap frontier are then suitable candidates for subsequent kernel allocations. This is important because the next phase forks thousands of child processes, forcing the kernel to allocate many \texttt{struct cred} objects close to the OOB-accessible region.

A relevant detail is the use of \texttt{MADV\_DONTFORK}. The exploit marks the large memory-pressure mapping so that child processes do not inherit it. Without this, forking thousands of children would duplicate large page-table structures and could exhaust memory before the credential spray phase.
\subsubsection{Credential Spray}

After calibration, the exploit \textbf{sprays credential objects} by forking thousands of child processes. Each child forces the kernel to allocate a separate \texttt{struct cred}; the children are then blocked on a pipe so that their credentials \textbf{remain allocated and stable} while the parent scans memory.

The target structure contains several identity fields:

\begin{itemize}
    \item \texttt{uid}, \texttt{gid};
    \item \texttt{suid}, \texttt{sgid};
    \item \texttt{euid}, \texttt{egid};
    \item \texttt{fsuid}, \texttt{fsgid}.
\end{itemize}

For a normal user in the test environment, these fields contain the value 1000. For root, they contain zero. The objective is therefore to locate a live child credential object and overwrite these fields with zero.

\subsubsection{Credential Scan}

The parent process scans the OOB-accessible memory window looking for the 64-bit signature:

\begin{lstlisting}[caption={Credential signature searched by the exploit.}]
0x000003E8000003E8
\end{lstlisting}

This value corresponds to two adjacent 32-bit fields containing decimal 1000, namely \texttt{uid=1000} and \texttt{gid=1000}. The exploit does not trust a single match. It verifies additional adjacent fields to confirm that the pattern belongs to a real \texttt{struct cred} and not to unrelated memory.

\begin{lstlisting}[caption={Credential discovery output.}]
[+] *** CRED FOUND at OOB offset 0xc59f8 ***
[+] Verified: 3 consecutive uid_sig matches
\end{lstlisting}

The discovered offset identifies the beginning of the useful credential identity fields for one of the sprayed child processes.

\subsubsection{Credential Overwrite}

Privilege escalation is then performed with four 8-byte writes that \textbf{zero the UID/GID-related fields} of the selected child credential object. The first write also touches the adjacent usage field, because the detected signature starts at the \texttt{uid} field.

\begin{lstlisting}[language=C, caption={Credential overwrite strategy.}]
oob_write(cred_off - 4,  0);  /* usage + uid  */
oob_write(cred_off + 4,  0);  /* gid + suid   */
oob_write(cred_off + 12, 0);  /* sgid + euid  */
oob_write(cred_off + 20, 0);  /* egid + fsuid */
\end{lstlisting}

After the overwrite, the kernel reads UID/GID fields equal to zero and treats the target child process as root. This does not invoke \texttt{setuid()}, does not pass through the normal authorization logic and does not generate normal authentication logs. The kernel identity state is changed directly in memory.

A verification read confirms the operation:

\begin{lstlisting}[caption={Credential overwrite verification.}]
[*] Verify (usage+uid): 0x0000000000000000 -> ZEROED!
\end{lstlisting}

\subsubsection{Root Shell Creation}

Once the parent closes the synchronization pipe, all children wake up and call \texttt{getuid()}. Only the child whose credentials were overwritten observes UID 0. That child then copies the exploit binary to \texttt{/tmp/rootsh}, changes ownership to root and enables the SUID bit.

The exploit binary also contains a SUID entry path:

\begin{lstlisting}[language=C, caption={SUID payload entry.}]
if (geteuid() == 0) {
    setuid(0);
    setgid(0);
    execl("/bin/sh", "sh", NULL);
}
\end{lstlisting}

After the SUID copy is created, executing \texttt{/tmp/rootsh} yields an interactive root shell:

\begin{lstlisting}[caption={Final result.}]
$ ./rootsh
[+] SUID payload triggered! Spawning root shell...
$ whoami
root
\end{lstlisting}

\subsection{Experimental Results}

The execution of the target exploit produced the following log output:

\begin{figure}[H]
    \centering
    \includegraphics[width=0.8\textwidth]{img/CVE-2023-39191.png}
    \caption{Exploit execution log for CVE-2023-39191.}
    \label{fig:log-2023-39191}
\end{figure}

The full exploit chain was successfully executed in the vulnerable test environment. The important result is that the exploit did not merely demonstrate an OOB read: it reached a complete LPE by modifying a live \texttt{struct cred} and producing a persistent SUID-root helper for the lifetime of the filesystem.

The observed exploitation path was:

\begin{figure}[H]
\centering
\begin{tikzpicture}[
    node distance=0.55cm,
    every node/.style={
        draw,
        rounded corners,
        align=center,
        font=\scriptsize,
        minimum height=0.8cm,
        text width=2.25cm
    },
    arrow/.style={-{Latex[length=2mm]}, thick}
]

\node (v) {Verifier accepts\\overlap};
\node[right=of v] (rw) {OOB R/W\\primitive};
\node[right=of rw] (cal) {Memory\\calibration};
\node[right=of cal] (spray) {Credential\\spray};

\node[below=0.75cm of spray] (scan) {UID/GID\\scan};
\node[left=of scan] (zero) {Credential\\overwrite};
\node[left=of zero] (suid) {Create\\\texttt{/tmp/rootsh}};
\node[left=of suid] (root) {Root\\shell};

\draw[arrow] (v) -- (rw);
\draw[arrow] (rw) -- (cal);
\draw[arrow] (cal) -- (spray);
\draw[arrow] (spray) -- (scan);
\draw[arrow] (scan) -- (zero);
\draw[arrow] (zero) -- (suid);
\draw[arrow] (suid) -- (root);

\end{tikzpicture}
\caption{CVE-2023-39191 exploitation path.}
\end{figure}

From a stability perspective, the exploit is notable because it modifies a single credential object rather than global kernel control structures. The system remains operational after exploitation. The modified credential object is limited to the affected child process, while the SUID-root payload remains available until reboot or filesystem cleanup.

This makes the exploit stealthier than crash-oriented or global-hooking kernel attacks. There is no legitimate authentication event, because privilege escalation occurs by directly modifying kernel memory. The most visible artifact is the unexpected SUID-root binary created under \texttt{/tmp}.

\subsection{Comparison with Previous eBPF Exploits}

Among the exploits analyzed in the previous reports, the closest comparison for CVE-2023-39191 is CVE-2023-2163. The concrete verifier bug is different, but the exploitation model is very similar: both attacks exploit a mismatch between the verifier's abstract interpretation and the actual runtime state.

In CVE-2023-2163, incorrect path pruning causes the verifier to skip a branch where pointer bounds have changed. Real execution can still take that branch, producing a one-byte stack overflow. The exploit then brute-forces the corrupted byte, leaks a map pointer, upgrades the primitive to arbitrary 64-bit read/write, and finally patches credentials to obtain root. In CVE-2023-39191, the stale verifier state is not produced by path pruning, but by incomplete dynptr invalidation: the verifier still trusts a dynptr whose metadata has been partially overwritten at runtime.

\begin{table}[H]
\centering
\begin{tabularx}{\textwidth}{|l|X|X|}
\hline
\textbf{Aspect} & \textbf{CVE-2023-2163} & \textbf{CVE-2023-39191} \\
\hline
Verifier mistake & Incorrect path pruning & Missing dynptr invalidation after partial stack overlap \\
\hline
Divergence & Pruned path differs from runtime path & Tracked dynptr differs from runtime object \\
\hline
Initial primitive & One-byte stack overflow & Corrupted dynptr metadata \\
\hline
Primitive upgrade & Brute-force low byte, leak map pointer, build arbitrary R/W & Corrupt dynptr size/offset, calibrate memory, build OOB R/W \\
\hline
Exploit shaping & Pointer leak, brute force, arbitrary R/W construction & Adaptive memory calibration, cred spray, OOB scan \\
\hline
Final step & Patch \texttt{cred} and spawn root shell & Spray/scan \texttt{cred}, zero UID/GID, spawn root shell \\
\hline
\end{tabularx}
\caption{Exploit-level parallelism between CVE-2023-2163 and CVE-2023-39191.}
\end{table}

\subsection{Patch and Mitigation}

The vulnerability was fixed by strengthening how the verifier handles dynptr stack slots and stack-slot reuse. The relevant patch series improves dynptr stack-state invalidation, including invalidation of slices when dynptrs on the stack are destroyed and stricter handling of dynptr stack-slot reinitialization. Conceptually, whenever a helper writes to stack memory that overlaps a tracked dynptr, the previous dynptr state must no longer remain trusted.

\begin{lstlisting}[caption={Conceptual model of the fix.}]
if stack write overlaps tracked dynptr state:
    invalidate previous dynptr state
    mark affected stack slots/slices as unsafe
\end{lstlisting}

This closes the verifier/runtime gap because the first \texttt{dynptr} can no longer survive a partial overwrite of its metadata. After the fix, the overlap pattern is rejected or the stale object is invalidated before it can be used for a corrupted dynptr access.

The most relevant mitigations are:

\begin{itemize}
    \item upgrade to a kernel containing the dynptr invalidation fix;
    \item disable unprivileged BPF where possible;
    \item restrict BPF program loading to trusted users;
    \item monitor suspicious SUID file creation such as \texttt{/tmp/rootsh};
    \item monitor unusual BPF map/ring-buffer activity by untrusted users.
\end{itemize}

\subsection{Security Impact}

CVE-2023-39191 demonstrates why verifier correctness is a critical kernel security boundary. The BPF runtime trusts the verifier's abstract model: once a program is accepted, the kernel does not re-check every high-level safety property during execution. If the abstract model diverges from the actual runtime object layout, the attacker obtains kernel memory access while still executing verifier-approved code.

The impact is especially severe because the exploit reaches full LPE without corrupting global kernel structures. By modifying a single live \texttt{struct cred}, the attack changes the identity of one child process while leaving the rest of the system stable. This lowers the chance of immediate crashes and makes exploitation more reliable.

The defensive lesson is that object lifetime and object integrity must be tracked at byte granularity when the verifier models compound stack objects. A \texttt{dynptr} is not safe if only part of it remains valid. Partial overwrite must be treated exactly like full destruction, because metadata corruption is enough to transform a bounded object into a large OOB access window.

In short, the verifier approved what it believed to be a safe local dynptr, but the kernel executed operations on a corrupted runtime object with an approximately 4 GiB access window in the tested exploit configuration. That mismatch is sufficient to move from a local BPF program to root privileges.

\end{document}